% --- Archon physics formalization source begin ---
% archon:physics
% archon:covers IPhO2026Problems/problem_IPhO_2026_4_B_6.lean
% archon:source-report reports/ipho_2026/problem_IPhO_2026_4_B_6.source.json
% archon:problem-id IPhO_2026_4
% archon:part-id B.6
% archon:previous-part IPhO_2026_4_B_5
% archon:previous-part-policy natural_language_prerequisite_only

\chapter{Physics problem IPhO\_2026\_4\_B\_6}
\label{ch:IPhO2026Problems_problem_IPhO_2026_4_B_6}

\paragraph{Problem source.}
The inner cylinder contains dry air plus water vapor at total pressure
approximately P\_atm.  The water level is adjusted and its height H is recorded
as temperature T falls.  At T\_0 = 273.15 K, extrapolated height is H\_0 and the
water vapor pressure may be taken as zero.  Vapor pressure obeys
ln(P\_v/P\_v0) = -(Q\_v/R)*(1/T - 1/T\_0).  Use the experimental procedure and
geometry on pages 11--12.

Current subquestion:
Convert Q\_v into latent heat per unit mass L\_v and state the formula.

\paragraph{Current subquestion.}
Convert Q\_v into latent heat per unit mass L\_v and state the formula.

\paragraph{Recorded answer/context.}
L\_v = Q\_v/M\_0 = 2190 +/- 110 kJ/kg.

\paragraph{Figure/image path.}
/root/proposal\_for\_physic/science-mango/ipho\_2026\_source/image/E1\_page-12.png

\paragraph{Reusable previous-part conclusions.}
\begin{itemize}
\item Source B.5. Question: Construct a Clausius-Clapeyron graph and use it to determine the molar latent heat Q\_v. Reusable conclusions: Plot ln(P\_v/P\_atm) against 1/T; official sample slope is -4700 +/- 200 K and Q\_v = 39 +/- 2 kJ/mol. Policy: natural\_language\_prerequisite\_only; do\_not\_import\_Lean\_output
\end{itemize}

\paragraph{Formalization target.}
create a compiling Lean file with sorry bodies at `IPhO2026Problems/problem\_IPhO\_2026\_4\_B\_6.lean`. The Lean declarations must preserve the physical quantities, dimensions or dimensional roles, figure labels, governing-law hypotheses, and final relation expressed by this problem.
Use Mathlib/Physlib names found through LeanExplore where available. If a physics API is missing, introduce faithful local abstractions rather than scalar placeholder aliases.

\begin{definition}[energyDimension]
  \label{decl:physics:IPhO_2026_4_B_6:energyDimension}
  \lean{IPhO2026Problems.IPhO2026_4_B_6.energyDimension}
  The mechanical dimension of energy.
\end{definition}
\begin{proof}
  This is the defining declaration; unfolding it gives the stated typed data or relation.
\end{proof}

\begin{definition}[specificEnergyDimension]
  \label{decl:physics:IPhO_2026_4_B_6:specificEnergyDimension}
  \lean{IPhO2026Problems.IPhO2026_4_B_6.specificEnergyDimension}
  \uses{decl:physics:IPhO_2026_4_B_6:energyDimension}
  The dimension of energy per unit mass.
\end{definition}
\begin{proof}
  This is the defining declaration; unfolding it gives the stated typed data or relation.
\end{proof}

\begin{definition}[Temperature]
  \label{decl:physics:IPhO_2026_4_B_6:Temperature}
  \lean{IPhO2026Problems.IPhO2026_4_B_6.Temperature}
  The local abbreviation Temperature records part of the typed physics model.
\end{definition}
\begin{proof}
  This is the defining declaration; unfolding it gives the stated typed data or relation.
\end{proof}

\begin{definition}[Length]
  \label{decl:physics:IPhO_2026_4_B_6:Length}
  \lean{IPhO2026Problems.IPhO2026_4_B_6.Length}
  The local abbreviation Length records part of the typed physics model.
\end{definition}
\begin{proof}
  This is the defining declaration; unfolding it gives the stated typed data or relation.
\end{proof}

\begin{definition}[Pressure]
  \label{decl:physics:IPhO_2026_4_B_6:Pressure}
  \lean{IPhO2026Problems.IPhO2026_4_B_6.Pressure}
  The local abbreviation Pressure records part of the typed physics model.
\end{definition}
\begin{proof}
  This is the defining declaration; unfolding it gives the stated typed data or relation.
\end{proof}

\begin{definition}[Mass]
  \label{decl:physics:IPhO_2026_4_B_6:Mass}
  \lean{IPhO2026Problems.IPhO2026_4_B_6.Mass}
  The local abbreviation Mass records part of the typed physics model.
\end{definition}
\begin{proof}
  This is the defining declaration; unfolding it gives the stated typed data or relation.
\end{proof}

\begin{definition}[Energy]
  \label{decl:physics:IPhO_2026_4_B_6:Energy}
  \lean{IPhO2026Problems.IPhO2026_4_B_6.Energy}
  \uses{decl:physics:IPhO_2026_4_B_6:energyDimension}
  The local abbreviation Energy records part of the typed physics model.
\end{definition}
\begin{proof}
  This is the defining declaration; unfolding it gives the stated typed data or relation.
\end{proof}

\begin{definition}[SpecificEnergy]
  \label{decl:physics:IPhO_2026_4_B_6:SpecificEnergy}
  \lean{IPhO2026Problems.IPhO2026_4_B_6.SpecificEnergy}
  \uses{decl:physics:IPhO_2026_4_B_6:specificEnergyDimension}
  The local abbreviation SpecificEnergy records part of the typed physics model.
\end{definition}
\begin{proof}
  This is the defining declaration; unfolding it gives the stated typed data or relation.
\end{proof}

\begin{definition}[siReadout]
  \label{decl:physics:IPhO_2026_4_B_6:siReadout}
  \lean{IPhO2026Problems.IPhO2026_4_B_6.siReadout}
  The scalar readout of a dimensionful quantity in the SI unit system.
\end{definition}
\begin{proof}
  This is the defining declaration; unfolding it gives the stated typed data or relation.
\end{proof}

\begin{definition}[MolarLatentHeatEstimate]
  \label{decl:physics:IPhO_2026_4_B_6:MolarLatentHeatEstimate}
  \lean{IPhO2026Problems.IPhO2026_4_B_6.MolarLatentHeatEstimate}
  The part B.5 estimate of molar latent heat.
\end{definition}
\begin{proof}
  This is the defining declaration; unfolding it gives the stated typed data or relation.
\end{proof}

\begin{definition}[VaporizationExperiment]
  \label{decl:physics:IPhO_2026_4_B_6:VaporizationExperiment}
  \lean{IPhO2026Problems.IPhO2026_4_B_6.VaporizationExperiment}
  \uses{decl:physics:IPhO_2026_4_B_6:Temperature, decl:physics:IPhO_2026_4_B_6:Length, decl:physics:IPhO_2026_4_B_6:Pressure, decl:physics:IPhO_2026_4_B_6:Mass, decl:physics:IPhO_2026_4_B_6:Energy, decl:physics:IPhO_2026_4_B_6:SpecificEnergy, decl:physics:IPhO_2026_4_B_6:MolarLatentHeatEstimate}
  Quantities and labelled functions appearing in the B.1--B.6 inner-cylinder experiment.
\end{definition}
\begin{proof}
  This is the defining declaration; unfolding it gives the stated typed data or relation.
\end{proof}

\begin{definition}[HasReferenceAndProcedureData]
  \label{decl:physics:IPhO_2026_4_B_6:HasReferenceAndProcedureData}
  \lean{IPhO2026Problems.IPhO2026_4_B_6.HasReferenceAndProcedureData}
  \uses{decl:physics:IPhO_2026_4_B_6:Temperature, decl:physics:IPhO_2026_4_B_6:siReadout, decl:physics:IPhO_2026_4_B_6:VaporizationExperiment}
  Reference values and procedure readouts from pages 11--12, together with the standard molar mass of water needed for B.6.
\end{definition}
\begin{proof}
  This is the defining declaration; unfolding it gives the stated typed data or relation.
\end{proof}

\begin{definition}[GoverningLaws]
  \label{decl:physics:IPhO_2026_4_B_6:GoverningLaws}
  \lean{IPhO2026Problems.IPhO2026_4_B_6.GoverningLaws}
  \uses{decl:physics:IPhO_2026_4_B_6:Temperature, decl:physics:IPhO_2026_4_B_6:siReadout, decl:physics:IPhO_2026_4_B_6:VaporizationExperiment}
  The physical laws used by the experimental model.
\end{definition}
\begin{proof}
  This is the defining declaration; unfolding it gives the stated typed data or relation.
\end{proof}

\begin{definition}[PreviousPartB5Result]
  \label{decl:physics:IPhO_2026_4_B_6:PreviousPartB5Result}
  \lean{IPhO2026Problems.IPhO2026_4_B_6.PreviousPartB5Result}
  \uses{decl:physics:IPhO_2026_4_B_6:VaporizationExperiment}
  The only imported previous-part conclusion: the B.5 graph has reported slope “-4700 ± 200 K” and gives “Q\_v = 39 ± 2 kJ/mol”.
\end{definition}
\begin{proof}
  This is the defining declaration; unfolding it gives the stated typed data or relation.
\end{proof}

\begin{theorem}[Physics formalization target]
\label{thm:physics:IPhO_2026_4_B_6:target}
\lean{IPhO2026Problems.IPhO2026_4_B_6.latentHeatPerUnitMass_from_molarEstimate}
\uses{decl:physics:IPhO_2026_4_B_6:energyDimension, decl:physics:IPhO_2026_4_B_6:specificEnergyDimension, decl:physics:IPhO_2026_4_B_6:Temperature, decl:physics:IPhO_2026_4_B_6:Length, decl:physics:IPhO_2026_4_B_6:Pressure, decl:physics:IPhO_2026_4_B_6:Mass, decl:physics:IPhO_2026_4_B_6:Energy, decl:physics:IPhO_2026_4_B_6:SpecificEnergy, decl:physics:IPhO_2026_4_B_6:siReadout, decl:physics:IPhO_2026_4_B_6:MolarLatentHeatEstimate, decl:physics:IPhO_2026_4_B_6:VaporizationExperiment, decl:physics:IPhO_2026_4_B_6:HasReferenceAndProcedureData, decl:physics:IPhO_2026_4_B_6:GoverningLaws, decl:physics:IPhO_2026_4_B_6:PreviousPartB5Result}
The assigned autoformalize agent should translate this physics problem into Lean declarations in the covered file, with theorem and lemma proof bodies written as `by sorry`.
\end{theorem}
\begin{proof}
This is an autoformalization task, not a proof task. Produce statements that can later be proved by the physics prover without weakening the physical model.
\end{proof}
% --- Archon physics formalization source end ---
